\section{Large-Budget Theory for Weighted Target-Impact Objectives}
\label{sec:large-budget-general-prior-appendix}

We now turn to the large-budget behavior of the paper's main objective $V=V_{\twoimpact}$.
As with the small-budget results, we provide proofs in a more general form and then specialize it to the named objectives in the paper.
The relevant large-budget class consists of weighted target-impact objectives
\[
  V_Z(\bm s_1)
  \defeq
  \E[\Test_1\cdot Z\cdot \ATE_3],
\]
where $Z$ is a downstream non-negative weight that is non-adaptive with respect to the pilot-stage sampling noise, meaning $Z\ind \varepsilon_1\mid\bm\tau$.
The paper's main objective $V_\twoimpact$ takes this form for $Z=\Test_2$, and the single-stage objective $V_\singleimpact$ takes this form for $Z=1$.

Throughout this section, we will use the unit-investment design set
\[
  \overline{\mathcal S}
  =
  \{\overline{\bm s}_1\in\R_{\ge0}^D:\bm c^\top\overline{\bm s}_1=1\}
\]
and the representative unit-investment design
\[
  \overline{\bm s}_1^{\,\infty}
  \defeq
  \frac{\bm s_3}{\bm c^\top\bm s_3}.
\]

We will rely on three prior conditions for the proofs.
First, $\ATE_3$ has a finite second moment and is \emph{non-degenerate}, i.e. has non-zero variance.
Second, the prior satisfies what we call a \emph{uniform boundary anti-concentration} condition.
With
\[
  \xi(\overline{\bm s}_1,\bm\tau)
  \defeq
  \frac{\overline{\bm s}_1^\top\bm\tau}
  {\sqrt{\bm1^\top\overline{\bm s}_1}},
\]
we assume
\begin{equation}
  \omega(\delta)
  \defeq
  \sup_{\overline{\bm s}_1\in\overline{\mathcal S}}
  \P\!\left(|\xi(\overline{\bm s}_1,\bm\tau)|\le \delta\right)
  \xrightarrow[\delta\downarrow0]{}0.
  \label{eq:uniform-boundary-anti-concentration-large-budget-appendix}
\end{equation}
Third, the prior satisfies the \emph{sign-separation condition}
\begin{equation}
  \forall \overline{\bm s}_1\in\overline{\mathcal{S}}\text{ not proportional to }\bm s_3:\qquad
  \P\!\big((\overline{\bm s}_1^\top\bm\tau)(\bm s_3^\top\bm\tau)<0\big)>0.
  \label{eq:ordinary-sign-separation-large-budget-appendix}
\end{equation}
Distributions with a Lebesgue density positive on an open ball containing the origin satisfy the latter two regularity conditions, and they also imply positive variance of $\ATE_3$ whenever $\E[\ATE_3^2]<\infty$; see \Cref{prop:density-large-budget-regularity-appendix}.

For any non-adaptive downstream non-negative weight $Z$, write
\[
  m_Z(\bm\tau)
  \defeq
  \E[Z\mid\bm\tau].
\]
We define the corresponding follow-up, oracle, and noise-free values as
\begin{align*}
  V_{Z,\followup}
  &\defeq
  \E[Z\cdot \ATE_3],\\
  V_{Z,\oracle}
  &\defeq
  \E[\indic{\ATE_3>0}\cdot Z\cdot \ATE_3],\\
  V_{Z,\noisefree}(\overline{\bm s}_1)
  &\defeq
  \E[\indic{\overline{\bm s}_1^\top\bm\tau>0}\cdot Z\cdot \ATE_3].
\end{align*}
For a budget $b>0$, define the budget-constrained value and solution set
\begin{align*}
  V_Z^\star(b)
  &\defeq
  \sup_{\substack{\bm s_1\in\R_{\geq 0}^{D} \setminus \{\bm{0}\}\\ \bm c^\top\bm s_1\le b}}
  V_Z(\bm s_1),\\
  \mathcal S_{1,Z}^\star(b)
  &\defeq
  \arg\max_{\substack{\bm s_1\in\R_{\geq 0}^{D} \setminus \{\bm{0}\}\\ \bm c^\top\bm s_1\le b}}
  V_Z(\bm s_1).
\end{align*}
When $Z=\Test_2$, these reduce to the paper's main value $V^\star(b)$ and solution set $\mathcal S_1^\star(b)$.

The weighted objective requires the following two analogues of nondegeneracy and sign separation.
First, we assume \emph{weighted nondegeneracy}:
\begin{equation}
  \E[m_Z(\bm\tau)|\ATE_3|]>0.
  \label{eq:weighted-nondegeneracy}
\end{equation}
Second, we assume \emph{weighted sign separation}:
\begin{equation}
  \forall \overline{\bm s}_1\in\overline{\mathcal{S}}\text{ not proportional to }\bm s_3:\qquad
  \E\!\left[
    m_Z(\bm\tau)|\ATE_3|
    \indic{(\overline{\bm s}_1^\top\bm\tau)(\bm s_3^\top\bm\tau)<0}
  \right]>0.
  \label{eq:weighted-sign-separation}
\end{equation}
These conditions are automatic for the paper's objective once the ordinary sign-separation and non-zero-variance conditions hold, because $\E[\Test_2\mid\bm\tau]=p_2(\ATE_2)>0$ almost surely; this verification is given in \Cref{lem:test-two-satisfies-weighted-conditions-appendix}.

The main generic result of this section is the following theorem.
The rest of the section proves it and then specializes it to $Z=\Test_2$.
The proof relies on three key facts: finite-investment pilots are strictly below the oracle value, large-investment pilots converge uniformly to a noise-free pilot rule, and the unique noise-free design that attains the oracle value is the representative design.

\begin{theorem}[Large-budget limit for weighted target-impact objectives]
  \label{thm:large-budget-weighted-target-impact-appendix}
  Let $Z$ be non-adaptive and suppose $0\le Z\le M<\infty$ almost surely.
  Suppose $\E[\ATE_3^2]<\infty$, the uniform boundary anti-concentration condition \Cref{eq:uniform-boundary-anti-concentration-large-budget-appendix} holds, and the weighted non-degeneracy and weighted sign-separation conditions \Cref{eq:weighted-nondegeneracy,eq:weighted-sign-separation} hold.
  Then $\mathcal S_{1,Z}^\star(b)\neq\emptyset$ for all $b>0$ and:
  \begin{enumerate}
    \item[\rm (i)] $V_Z^\star(b)\le V_{Z,\oracle}$ for all $b>0$.
    \item[\rm (ii)] $V_Z^\star(b)\to V_{Z,\oracle}$ as $b\to\infty$.
    \item[\rm (iii)] The optimal investment diverges,
      \[
        \inf_{\bm s_1^\star\in\mathcal S_{1,Z}^\star(b)} \bm c^\top\bm s_1^\star
        \to\infty
        \qquad\text{as }b\to\infty.
      \]
    \item[\rm (iv)] The set of unit-investment normalized optimal designs converges to the unit-investment representative design,
      \[
        \sup_{\bm s_1^\star\in\mathcal S_{1,Z}^\star(b)}
        \left\|
        \frac{\bm s_1^\star}{\bm c^\top\bm s_1^\star}
        -
        \overline{\bm s}_1^{\,\infty}
        \right\|
        \to0
        \qquad\text{as }b\to\infty.
      \]
  \end{enumerate}
\end{theorem}

\subsection*{Proof ingredients for the generic theorem}

\begin{lemma}[Boundary and noise-free limits]
  \label{prop:limiting-impact-values}
  Let $Z$ be non-adaptive and suppose $\E[|Z \cdot \ATE_3|]<\infty$.
  Then $V_Z$ is continuous on $\R_{\geq 0}^{D} \setminus \{\bm{0}\}$ and
  \[
    \lim_{\bm s_1\to\bm0}V_Z(\bm s_1)
    =
    \alpha_1\,V_{Z,\followup}.
  \]
  If, in addition, the uniform boundary anti-concentration condition \Cref{eq:uniform-boundary-anti-concentration-large-budget-appendix} holds, then for every $\overline{\bm s}_1\in\overline{\mathcal{S}}$,
  \[
    \lim_{\invst\to\infty}V_Z(\invst\,\overline{\bm s}_1)
    =
    V_{Z,\noisefree}(\overline{\bm s}_1).
  \]
  Under the same condition, $V_{Z,\noisefree}$ is continuous on $\overline{\mathcal{S}}$.
\end{lemma}

\begin{proof}
  \emph{Step 1: Continuity of $V_Z$ on $\R_{\geq 0}^{D} \setminus \{\bm{0}\}$ and the small-sample limit.}
  For $\bm s_1\in\R_{\geq 0}^{D} \setminus \{\bm{0}\}$, write $n_1=\bm1^\top\bm s_1$ and
  \[
    q(\bm s_1,\bm\tau)
    \defeq
    \Phi\!\left(
      \frac{\bm s_1^\top\bm\tau}{\sqrt{n_1}}
      -
      z_{1-\alpha_1}
    \right).
  \]
  Since $q(\bm s_1,\bm\tau)=\P(\Test_1=1\mid\bm\tau)$ and $Z$ is independent of the pilot-stage sampling noise conditional on $\bm\tau$,
  \[
    V_Z(\bm s_1)
    =
    \E[q(\bm s_1,\bm\tau)\cdot Z\cdot \ATE_3].
  \]
  Fix any sequence $\bm s_1^{(k)}\to \bm s_1$ in $\R_{\geq 0}^{D} \setminus \{\bm{0}\}$.
  Then $q(\bm s_1^{(k)},\bm\tau)\to q(\bm s_1,\bm\tau)$ pointwise in $\bm\tau$ by continuity of the Gaussian distribution function.
  Since $0\le q(\bm s_1^{(k)},\bm\tau)\le 1$ and $|Z \cdot \ATE_3|$ is integrable, dominated convergence gives
  \[
    V_Z(\bm s_1^{(k)})\to V_Z(\bm s_1),
  \]
  proving continuity on $\R_{\geq 0}^{D} \setminus \{\bm{0}\}$.

  If $\bm s_1\to\bm0$, then $n_1\to0$ and
  \[
    \left|\frac{\bm s_1^\top\bm\tau}{\sqrt{n_1}}\right|
    \le
    \sqrt{n_1}\,\|\bm\tau\|_1
    \to 0
    \qquad\text{for each fixed }\bm\tau,
  \]
  so $q(\bm s_1,\bm\tau)\to\Phi(-z_{1-\alpha_1})=\alpha_1$ pointwise.
  Dominated convergence therefore gives
  \[
    \lim_{\bm s_1\to\bm0}V_Z(\bm s_1)
    =
    \alpha_1\,\E[Z\cdot \ATE_3]
    =
    \alpha_1\,V_{Z,\followup}.
  \]

  \emph{Step 2: Pointwise large-investment limit.}
  Fix $\overline{\bm s}_1\in\overline{\mathcal{S}}$ and write
  \[
    q_{\invst}(\bm\tau)
    \defeq
    \Phi\!\left(
      \sqrt{\invst}\,\xi(\overline{\bm s}_1,\bm\tau)
      -
      z_{1-\alpha_1}
    \right).
  \]
  The uniform boundary anti-concentration condition implies
  \[
    \P\!\big(
      \xi(\overline{\bm s}_1,\bm\tau)=0
    \big)=0.
  \]
  Hence $q_{\invst}(\bm\tau)\to \indic{\xi(\overline{\bm s}_1,\bm\tau)>0}$ almost surely as $\invst\to\infty$.
  Since $\xi(\overline{\bm s}_1,\bm\tau)>0$ if and only if $\overline{\bm s}_1^\top\bm\tau>0$, dominated convergence gives
  \[
    \lim_{\invst\to\infty}V_Z(\invst\,\overline{\bm s}_1)
    =
    \E[\indic{\overline{\bm s}_1^\top\bm\tau>0}\cdot Z\cdot \ATE_3]
    =
    V_{Z,\noisefree}(\overline{\bm s}_1).
  \]

  \emph{Step 3: Continuity of $V_{Z,\noisefree}$ on $\overline{\mathcal{S}}$.}
  Let $\overline{\bm s}_1^{(k)}\to\overline{\bm s}_1$ in $\overline{\mathcal{S}}$.
  Then $\xi(\overline{\bm s}_1^{(k)},\bm\tau)\to\xi(\overline{\bm s}_1,\bm\tau)$ pointwise in $\bm\tau$.
  By the previous argument,
  \[
    \P\!\big(\xi(\overline{\bm s}_1,\bm\tau)=0\big)=0,
  \]
  so
  \[
    \indic{\xi(\overline{\bm s}_1^{(k)},\bm\tau)>0}
    \to
    \indic{\xi(\overline{\bm s}_1,\bm\tau)>0}
    \qquad\text{almost surely.}
  \]
  Since $|\indic{\xi>0}\cdot Z\cdot \ATE_3|\le |Z \cdot \ATE_3|$, dominated convergence gives
  \[
    V_{Z,\noisefree}(\overline{\bm s}_1^{(k)})
    \to
    V_{Z,\noisefree}(\overline{\bm s}_1),
  \]
  proving continuity on $\overline{\mathcal{S}}$.
\end{proof}

\begin{proposition}[Uniform convergence to the noise-free pilot]
  \label{prop:large-investment-expansion-weighted-appendix}
  Let $Z$ be non-adaptive and suppose $0\le Z\le M<\infty$ almost surely.
  Suppose $\E[\ATE_3^2]<\infty$ and the uniform boundary anti-concentration condition \Cref{eq:uniform-boundary-anti-concentration-large-budget-appendix} holds.
  Then as $\invst\to\infty$,
  \[
    \sup_{\overline{\bm s}_1\in\overline{\mathcal{S}}}
    \big|V_Z(\invst\,\overline{\bm s}_1)-V_{Z,\noisefree}(\overline{\bm s}_1)\big|
    \to 0.
  \]
\end{proposition}

\begin{proof}
  Fix $\delta>0$ and bound the tail error by
  \[
    r_{\invst}(\delta)
    \defeq
    \max\!\left\{
      \Phi(z_{1-\alpha_1}-\sqrt{\invst}\,\delta),
      \Phi(-z_{1-\alpha_1}-\sqrt{\invst}\,\delta)
    \right\}.
  \]
  For every $\overline{\bm s}_1\in\overline{\mathcal{S}}$ and every realization of $\bm\tau$, bound the approximation error between the finite and infinite investment decisions
  \[
    e_{\invst}(\overline{\bm s}_1,\bm\tau)
    \defeq
    \left|
    \Phi\!\left(
      \sqrt{\invst}\,\xi(\overline{\bm s}_1,\bm\tau)-z_{1-\alpha_1}
    \right)
    -
    \indic{\xi(\overline{\bm s}_1,\bm\tau)>0}
    \right|.
  \]
  If $|\xi(\overline{\bm s}_1,\bm\tau)|>\delta$, the error is bounded by $r_{\invst}(\delta)$; otherwise it is bounded by $1$.
  Hence
  \[
    e_{\invst}(\overline{\bm s}_1,\bm\tau)
    \le
    \indic{\big|\xi(\overline{\bm s}_1,\bm\tau)\big|\le\delta}
    +
    r_{\invst}(\delta).
  \]
  Therefore, applying non-adaptivity of $Z$,
  \begin{align*}
    \big|V_Z(\invst\,\overline{\bm s}_1)-V_{Z,\noisefree}(\overline{\bm s}_1)\big|
    &=
    \left|
    \E\!\left[
      Z\cdot \ATE_3\cdot
      \left(
        \Phi\!\left(
          \sqrt{\invst}\,\xi(\overline{\bm s}_1,\bm\tau)-z_{1-\alpha_1}
        \right)
        -
        \indic{\xi(\overline{\bm s}_1,\bm\tau)>0}
      \right)
    \right]
    \right| \\
    &\le
    \E\!\left[
      Z\cdot |\ATE_3|\cdot e_{\invst}(\overline{\bm s}_1,\bm\tau)
    \right] \\
    &\le
    \E\!\left[
      Z \cdot |\ATE_3|\cdot
      \indic{\big|\xi(\overline{\bm s}_1,\bm\tau)\big|\le\delta}
    \right]
    +
    r_{\invst}(\delta)\,\E[Z\cdot |\ATE_3|].
  \end{align*}
  Using $0\le Z\le M$ and Cauchy--Schwarz,
  \[
    \E\!\left[
      Z\cdot |\ATE_3|\cdot
      \indic{\big|\xi(\overline{\bm s}_1,\bm\tau)\big|\le\delta}
    \right]
    \le
    M\sqrt{\E[\ATE_3^2]}
    \P\!\big(\big|\xi(\overline{\bm s}_1,\bm\tau)\big|\le\delta\big)^{1/2}
    \le
    M\sqrt{\E[\ATE_3^2] \omega(\delta)}.
  \]
  Taking the supremum over $\overline{\bm s}_1\in\overline{\mathcal{S}}$ gives
  \[
    \sup_{\overline{\bm s}_1\in\overline{\mathcal{S}}}
    \big|V_Z(\invst\,\overline{\bm s}_1)-V_{Z,\noisefree}(\overline{\bm s}_1)\big|
    \le
    M\sqrt{\E[\ATE_3^2]\omega(\delta)}
    +
    r_{\invst}(\delta)\,M\E[|\ATE_3|].
  \]

  Let $\Delta_{\invst}$ denote the left-hand side of the preceding inequality, and fix any $\delta>0$.
  Since $\alpha_1\in(0,1)$, $z_{1-\alpha_1}$ is finite, and hence both $z_{1-\alpha_1}-\sqrt{\invst}\delta$ and $-z_{1-\alpha_1}-\sqrt{\invst}\delta$ tend to $-\infty$ as $\invst\to\infty$.
  Therefore $r_{\invst}(\delta)\to0$ for each fixed $\delta>0$.
  Taking $\limsup_{\invst\to\infty}$ in the preceding bound gives
  \[
    \limsup_{\invst\to\infty}\Delta_{\invst}
    \le
    M\sqrt{\E[\ATE_3^2]\omega(\delta)} .
  \]
  Since this holds for every $\delta>0$, we may let $\delta\downarrow0$; the uniform boundary anti-concentration assumption gives $\omega(\delta)\to0$, so $\limsup_{\invst\to\infty}\Delta_{\invst}\le0$.
  As $\Delta_{\invst}\ge0$, it follows that $\Delta_{\invst}\to0$.
  This proves uniform convergence.
\end{proof}

\begin{lemma}[Existence of optimal designs]
  \label{prop:optimal-solution-exists-weighted-appendix}
  Let $Z$ be non-adaptive and suppose $\E[|Z \cdot \ATE_3|]<\infty$, $m_Z(\bm\tau)\ge0$ almost surely, and \Cref{eq:weighted-nondegeneracy} holds.
  Then $\mathcal S_{1,Z}^\star(b)\neq\emptyset$ for every $b>0$.
\end{lemma}

\begin{proof}
  Fix $b>0$ and define the extended reparametrized objective
  \[
    f_Z(\invst,\overline{\bm s}_1)\defeq
    \begin{cases}
      V_Z(\invst\,\overline{\bm s}_1), & \invst>0,\\
      \alpha_1\,V_{Z,\followup}, & \invst=0,
    \end{cases}
    \qquad
    (\invst,\overline{\bm s}_1)\in[0,b]\times\overline{\mathcal{S}}.
  \]
  By \Cref{prop:limiting-impact-values}, $f_Z$ is continuous on the compact set $[0,b]\times\overline{\mathcal{S}}$, so Weierstrass' theorem implies that $f_Z$ attains its maximum there.

  It remains to show that no maximizer has $\invst^\star=0$.
  Let $\overline n_1^\infty\defeq \bm1^\top\overline{\bm s}_1^{\,\infty}>0$ and define
  \[
    q_\invst(x)
    \defeq
    \Phi\!\big(\sqrt{\invst\,\overline n_1^\infty}\,x-z_{1-\alpha_1}\big).
  \]
  Because the representative pilot satisfies $\ATE_1=\ATE_3$, we have
  \[
    f_Z(\invst,\overline{\bm s}_1^{\,\infty})-
    \alpha_1V_{Z,\followup}
    =
    \E\!\left[
      m_Z(\bm\tau)\,\ATE_3\big(q_\invst(\ATE_3)-\alpha_1\big)
    \right].
  \]
  Since $q_\invst(0)=\alpha_1$ and $q_\invst(\cdot)$ is strictly increasing, the product
  \[
    x\big(q_\invst(x)-\alpha_1\big)
  \]
  is strictly positive for every $x\neq0$.
  Since $m_Z(\bm\tau)\ge0$ almost surely and \Cref{eq:weighted-nondegeneracy} holds, it follows that
  \[
    f_Z(\invst,\overline{\bm s}_1^{\,\infty})>
    \alpha_1V_{Z,\followup}
    =
    f_Z(0,\overline{\bm s}_1^{\,\infty})
    \qquad
    \forall \invst>0.
  \]
  Hence the value attained at any positive investment along the representative design strictly improves on $\invst=0$.
  No maximizer of $f_Z$ can therefore satisfy $\invst^\star=0$.

  Therefore there exists a maximizer $(\invst^\star,\overline{\bm s}_1^\star)$ with $\invst^\star>0$, and then $\invst^\star\overline{\bm s}_1^\star$ is an optimal feasible solution.
\end{proof}

\begin{lemma}[Oracle upper bound]
  \label{lem:oracle-upper-bound-weighted}
  Suppose $Z\ge0$ almost surely and $\E[Z\cdot |\ATE_3|]<\infty$.
  Define the oracle random payoff and oracle value by
  \[
    Y_Z
    \defeq
    \indic{\ATE_3>0}\cdot Z\cdot \ATE_3,
    \qquad
    V_{Z,\oracle}
    \defeq
    \E[Y_Z].
  \]
  Define realized payoff random variables under the noisy and noise-free pilot decisions as
  \[
    X_Z(\bm s_1)\defeq \Test_1\cdot Z\cdot \ATE_3,
    \qquad
    X_{Z,\infty}(\overline{\bm s}_1)
    \defeq
    \indic{\overline{\bm s}_1^\top\bm\tau>0}\cdot Z\cdot \ATE_3.
  \]
  Then for every pilot design $\bm s_1\in\R_{\geq 0}^{D} \setminus \{\bm{0}\}$ and every $\overline{\bm s}_1\in\overline{\mathcal{S}}$,
  \[
    X_Z(\bm s_1)\le Y_Z
    \qquad\text{and}\qquad
    X_{Z,\infty}(\overline{\bm s}_1)\le Y_Z
    \qquad\text{almost surely.}
  \]
  Consequently,
  \[
    V_Z(\bm s_1)\le V_{Z,\oracle}
    \qquad\text{and}\qquad
    V_{Z,\noisefree}(\overline{\bm s}_1)\le V_{Z,\oracle}.
  \]
\end{lemma}

\begin{proof}
  If $\ATE_3>0$, then $Y_Z=Z\cdot \ATE_3\ge0$ and any pilot decision only multiplies this payoff by a number in $\{0,1\}$.
  If $\ATE_3\le0$, then $Y_Z=0$ and any decision to continue yields the non-positive payoff $Z\cdot \ATE_3\le0$.
  Thus both the noisy and noise-free pilot payoffs are bounded above by $Y_Z$ almost surely.
  Taking expectations proves the result.
\end{proof}

\begin{proposition}[Unique optimal noise-free design]
  \label{prop:optimal-noise-free-pilot-design-weighted-appendix}
  Suppose $Z\ge0$ almost surely and $\E[Z\cdot|\ATE_3|]<\infty$.
  Then the unit-investment representative pilot design $\overline{\bm s}_1^{\,\infty}$ attains the oracle bound:
  \[
    V_{Z,\noisefree}(\overline{\bm s}_1^{\,\infty})
    =
    V_{Z,\oracle}.
  \]
  If the weighted sign-separation condition \Cref{eq:weighted-sign-separation} holds, then $\overline{\bm s}_1^{\,\infty}$ is the unique maximizer of $V_{Z,\noisefree}$ over $\overline{\mathcal{S}}$.
\end{proposition}

\begin{proof}
  Under the representative pilot,
  \[
    \ATE_1
    =
    \frac{\left(\overline{\bm s}_1^{\,\infty}\right)^\top\bm\tau}
    {\bm1^\top\overline{\bm s}_1^{\,\infty}}
    =
    \frac{\bm s_3^\top\bm\tau}{\bm1^\top\bm s_3}
    =
    \ATE_3.
  \]
  Hence $\indic{\ATE_1>0}=\indic{\ATE_3>0}$ almost surely under $\overline{\bm s}_1^{\,\infty}$, so
  \[
    V_{Z,\noisefree}(\overline{\bm s}_1^{\,\infty})
    =
    \E[\indic{\ATE_3>0}\cdot Z\cdot \ATE_3]
    =
    V_{Z,\oracle}.
  \]

  Now let $\overline{\bm s}_1\in\overline{\mathcal{S}}$ be any noise-free design.
  The oracle gap can be written as
  \[
    V_{Z,\oracle}-V_{Z,\noisefree}(\overline{\bm s}_1)
    =
    \E\!\left[
      \Big(
        \indic{\ATE_3>0}
        -
        \indic{\overline{\bm s}_1^\top\bm\tau>0}
      \Big)\cdot Z \cdot \ATE_3
    \right],
  \]
  and the integrand is nonnegative almost surely by \Cref{lem:oracle-upper-bound-weighted}.
  On the strict sign-disagreement event
  \[
    A(\overline{\bm s}_1)
    \defeq
    \{(\overline{\bm s}_1^\top\bm\tau)(\bm s_3^\top\bm\tau)<0\},
  \]
  the integrand is exactly $Z\cdot|\ATE_3|$.
  Therefore, conditioning on $\bm\tau$,
  \[
    V_{Z,\oracle}-V_{Z,\noisefree}(\overline{\bm s}_1)
    \ge
    \E\!\left[
      m_Z(\bm\tau)\cdot |\ATE_3|\cdot \indic{A(\overline{\bm s}_1)}
    \right].
  \]
  If $\overline{\bm s}_1$ is not proportional to $\bm s_3$, the right-hand side is strictly positive by \Cref{eq:weighted-sign-separation}.
  Thus no non-proportional design attains the oracle bound.
  Since the representative design does attain the oracle bound, every maximizer must be proportional to $\bm s_3$.
  The unit-investment constraint $\bm c^\top\overline{\bm s}_1=1$ then identifies the unique maximizer as
  \[
    \overline{\bm s}_1
    =
    \frac{\bm s_3}{\bm c^\top\bm s_3}
    =
    \overline{\bm s}_1^{\,\infty}.
  \]
\end{proof}

\begin{lemma}[Strict finite-investment oracle gap]
  \label{lem:strict-finite-investment-oracle-gap-weighted}
  Let $Z$ be non-adaptive and suppose $Z\ge0$ almost surely, $\E[Z\cdot|\ATE_3|]<\infty$, and \Cref{eq:weighted-nondegeneracy} holds.
  Then every finite-investment pilot design satisfies
  \[
    V_Z(\bm s_1)<V_{Z,\oracle}
    \qquad
    \forall \bm s_1\in\R_{\geq 0}^{D}\setminus\{\bm0\}.
  \]
  Consequently, if $\mathcal S_{1,Z}^\star(b)\neq\emptyset$ for a finite budget $b$, then
  \[
    V_Z^\star(b)<V_{Z,\oracle}.
  \]
\end{lemma}

\begin{proof}
  Fix any finite-investment design $\bm s_1\in\R_{\geq 0}^{D}\setminus\{\bm0\}$.
  Conditional on $\bm\tau$, the pilot test is independent of $Z$ and has pass probability
  \[
    p_1(\bm\tau)
    \defeq
    \P(\Test_1=1\mid\bm\tau)
    \in(0,1).
  \]
  With
  \[
    X_Z\defeq \Test_1\cdot Z\cdot \ATE_3,
    \qquad
    Y_Z\defeq \indic{\ATE_3>0}\cdot Z\cdot \ATE_3,
  \]
  the conditional oracle gap is
  \[
    \E[Y_Z-X_Z\mid\bm\tau]
    =
    m_Z(\bm\tau)|\ATE_3|
    \Big(
      (1-p_1(\bm\tau))\indic{\ATE_3>0}
      +
      p_1(\bm\tau)\indic{\ATE_3<0}
    \Big).
  \]
  This quantity is nonnegative everywhere and strictly positive wherever $m_Z(\bm\tau)|\ATE_3|>0$.
  By \Cref{eq:weighted-nondegeneracy}, the latter event has positive weighted mass, so
  \[
    V_Z(\bm s_1)=\E[X_Z]<\E[Y_Z]=V_{Z,\oracle}.
  \]
  If $\mathcal S_{1,Z}^\star(b)$ is nonempty, an optimizer under budget $b$ is a finite-investment design, so the strict inequality also holds for $V_Z^\star(b)$.
\end{proof}

Finally, we will use the following standard result, which follows, for example, from \cite[Prop.~7.15 and Thm.~7.33]{RockafellarWe09}.

\begin{lemma}[Uniform convergence transfers unique maximizers]
  \label{lem:uniform-argmax-transfer-large-budget}
  Let $K\subset \R^D$ be compact, let $f:K\to\R$ be continuous with unique maximizer $x^\star$, and let $f_\lambda:K\to\R$ be functions indexed by $\lambda\to\infty$ such that
  \[
    \sup_{x\in K}|f_\lambda(x)-f(x)|\to0.
  \]
  Then for every $\epsilon>0$ and any (equivalent) norm $\|\cdot\|$, there exists $\lambda_\epsilon<\infty$ such that, for all $\lambda\ge\lambda_\epsilon$, every maximizer of $f_\lambda$ over $K$ lies in $\{x\in K:\|x-x^\star\|<\epsilon\}$.
\end{lemma}

We now have all the components necessary to proceed with a proof of the main large-budget result.

\begin{proof}[Proof of \Cref{thm:large-budget-weighted-target-impact-appendix}]
  Non-emptiness of $\mathcal S_{1,Z}^\star(b)$ follows from \Cref{prop:optimal-solution-exists-weighted-appendix}, since $0\le Z\le M$ and $\E[\ATE_3^2]<\infty$ imply $\E[|Z \cdot \ATE_3|]<\infty$.

  Part (i) follows immediately from the oracle upper bound in \Cref{lem:oracle-upper-bound-weighted}.
  For the matching lower bound in part (ii), \Cref{prop:optimal-noise-free-pilot-design-weighted-appendix} gives
  \[
    V_{Z,\noisefree}(\overline{\bm s}_1^{\,\infty})=V_{Z,\oracle}.
  \]
  Therefore
  \[
    V_Z^\star(b)
    \ge
    V_Z(b\,\overline{\bm s}_1^{\,\infty})
    \ge
    V_{Z,\noisefree}(\overline{\bm s}_1^{\,\infty})
    -
    \sup_{\overline{\bm s}_1\in\overline{\mathcal{S}}}
    \big|V_Z(b\,\overline{\bm s}_1)-V_{Z,\noisefree}(\overline{\bm s}_1)\big|.
  \]
  The supremum term converges to $0$ by \Cref{prop:large-investment-expansion-weighted-appendix}, while part (i) supplies the upper bound $V_Z^\star(b)\le V_{Z,\oracle}$.
  Hence $V_Z^\star(b)\to V_{Z,\oracle}$.

  To prove part (iii), fix any finite $\invst_0>0$.
  By \Cref{lem:strict-finite-investment-oracle-gap-weighted},
  \[
    \delta_{\invst_0}
    \defeq
    V_{Z,\oracle}-V_Z^\star(\invst_0)
    >0.
  \]
  By part (ii), choose $B_{\invst_0}$ such that $V_Z^\star(b)>V_Z^\star(\invst_0)$ whenever $b\ge B_{\invst_0}$.
  If an optimal design under such a budget had investment at most $\invst_0$, it would be feasible for the budget-$\invst_0$ problem and could not exceed $V_Z^\star(\invst_0)$.
  This contradicts $V_Z^\star(b)>V_Z^\star(\invst_0)$.
  Since $\invst_0$ was arbitrary, the optimal investment diverges.

  For part (iv), define, for each investment $\invst>0$,
  \[
    g_\invst(\overline{\bm s}_1)
    \defeq
    V_Z(\invst\,\overline{\bm s}_1),
    \qquad
    g_\infty(\overline{\bm s}_1)
    \defeq
    V_{Z,\noisefree}(\overline{\bm s}_1).
  \]
  By \Cref{prop:large-investment-expansion-weighted-appendix}, $g_\invst\to g_\infty$ uniformly on $\overline{\mathcal S}$.
  By \Cref{prop:optimal-noise-free-pilot-design-weighted-appendix}, $g_\infty$ has the unique maximizer $\overline{\bm s}_1^{\,\infty}$.
  Applying \Cref{lem:uniform-argmax-transfer-large-budget} with $K=\overline{\mathcal S}$ shows that every maximizer of $g_\invst$ is arbitrarily close to $\overline{\bm s}_1^{\,\infty}$ once $\invst$ is large enough, where we note that by equivalence of norms in $\R^D$ the $1$-norm can be replaced with any other norm.

  Finally, let $\bm s_1^\star\in\mathcal S_{1,Z}^\star(b)$ and write
  \[
    \invst^\star\defeq \bm c^\top\bm s_1^\star,
    \qquad
    \overline{\bm s}_1^\star
    \defeq
    \frac{\bm s_1^\star}{\bm c^\top\bm s_1^\star}.
  \]
  For its fixed investment $\invst^\star$, the normalized design $\overline{\bm s}_1^\star$ must maximize $g_{\invst^\star}$ over $\overline{\mathcal S}$; otherwise replacing it by a better unit-investment composition at the same investment would produce a feasible design with a larger value.
  Part (iii) implies that $\invst^\star\to\infty$ uniformly over all optimal designs as $b\to\infty$.
  Combining this investment divergence with the argmax-transfer conclusion proves
  \[
    \sup_{\bm s_1^\star\in\mathcal S_{1,Z}^\star(b)}
    \left\|
    \frac{\bm s_1^\star}{\bm c^\top\bm s_1^\star}
    -
    \overline{\bm s}_1^{\,\infty}
    \right\|
    \to0.
  \]
\end{proof}

\subsection*{Specialization to the two-stage impact objective}

We next verify that the generic weighted assumptions hold for the paper's main two-stage impact objective. This follows because strictly positive conditional weights transfer ordinary sign separation to weighted sign separation.

\begin{lemma}[Positive weights transfer sign separation]
  \label{prop:positive-conditional-weight-appendix}
  Let $Z$ be non-adaptive and suppose $m_Z(\bm\tau)\ge0$ almost surely and $\E[m_Z(\bm\tau)|\ATE_3|]<\infty$.
  If $m_Z(\bm\tau)>0$ almost surely and $\var(\ATE_3)>0$, then \Cref{eq:weighted-nondegeneracy} holds.
  If $m_Z(\bm\tau)>0$ almost surely and the ordinary sign-separation condition \Cref{eq:ordinary-sign-separation-large-budget-appendix} holds, then the weighted sign-separation condition \Cref{eq:weighted-sign-separation} holds.
\end{lemma}

\begin{proof}
  Since $m_Z(\bm\tau)>0$ almost surely and $\var(\ATE_3)>0$, we have $m_Z(\bm\tau)|\ATE_3|>0$ on a set of positive probability.
  This proves \Cref{eq:weighted-nondegeneracy}.

  Now fix any $\overline{\bm s}_1\in\overline{\mathcal{S}}$ not proportional to $\bm s_3$, and let
  \[
    A(\overline{\bm s}_1)
    \defeq
    \{(\overline{\bm s}_1^\top\bm\tau)(\bm s_3^\top\bm\tau)<0\}.
  \]
  By the ordinary sign-separation condition, $\P(A(\overline{\bm s}_1))>0$.
  On this event, $|\ATE_3|>0$, and by assumption $m_Z(\bm\tau)>0$ almost surely.
  Hence the nonnegative random variable
  \[
    m_Z(\bm\tau)|\ATE_3|\indic{A(\overline{\bm s}_1)}
  \]
  is strictly positive on a set of positive probability, so its expectation is strictly positive.
  This is exactly \Cref{eq:weighted-sign-separation}.
\end{proof}

\begin{lemma}[The follow-up test satisfies the weight assumptions]
  \label{lem:test-two-satisfies-weighted-conditions-appendix}
  Let $Z=\Test_2$.
  Then $Z$ is non-adaptive, $0\le Z\le1$ almost surely, and
  \[
    m_Z(\bm\tau)
    =
    \E[\Test_2\mid\bm\tau]
    =
    p_2(\ATE_2)
    \in(0,1)
    \qquad\text{almost surely.}
  \]
  Consequently, if $\var(\ATE_3)>0$ and the ordinary sign-separation condition \Cref{eq:ordinary-sign-separation-large-budget-appendix} holds, then \Cref{eq:weighted-nondegeneracy,eq:weighted-sign-separation} hold for $Z=\Test_2$.
\end{lemma}

\begin{proof}
  The follow-up test $\Test_2$ depends only on $\bm\tau$ and the follow-up-stage sampling noise $\varepsilon_2$.
  Since $\varepsilon_2$ is independent of the pilot-stage sampling noise $\varepsilon_1$ conditional on $\bm\tau$, $\Test_2$ is non-adaptive.
  The bound $0\le\Test_2\le1$ is immediate.
  Finally, by the definition of the stage-2 pass probability,
  \[
    \E[\Test_2\mid\bm\tau]
    =
    \P(\Test_2=1\mid\bm\tau)
    =
    p_2(\ATE_2),
  \]
  and $p_2(\ATE_2)\in(0,1)$ for every finite $\ATE_2$.
  The final claim follows from \Cref{prop:positive-conditional-weight-appendix}.
\end{proof}

\refstepcounter{proposition}
\begin{linkedresult}{proof}{prop:optimal-solution-exists}[Existence of optimal solution]
  \label{prop:optimal-solution-exists-appendix}
  Suppose $\E[|\ATE_3|]<\infty$ and $\var(\ATE_3)>0$.
  Then the pilot impact maximization problem defined in \Cref{sec:pilot-budget-constraint-and-objective} has at least one optimal solution, i.e.\ $\mathcal S_1^\star(b)\neq\emptyset$ for all $b>0$.
\end{linkedresult}

\begin{proof}
  Apply \Cref{prop:optimal-solution-exists-weighted-appendix} with $Z=\Test_2$.
  By \Cref{lem:test-two-satisfies-weighted-conditions-appendix}, $Z$ is non-adaptive, $m_Z(\bm\tau)=p_2(\ATE_2)>0$ almost surely, and $0\le Z\le1$.
  Since $\var(\ATE_3)>0$, \Cref{eq:weighted-nondegeneracy} holds by \Cref{prop:positive-conditional-weight-appendix}.
  Finally, $\E[|Z \cdot \ATE_3|]\le\E[|\ATE_3|]<\infty$.
\end{proof}

The next result is a slightly more general version of the large-budget theorem for the pilot impact.
That the density condition used in the body theorem implies the conditions here follows from \Cref{prop:density-large-budget-regularity-appendix} below.

\refstepcounter{theorem}
\begin{linkedresult}{extended}{thm:large-budget-limiting-optimal-policy}
  \label{thm:large-budget-limiting-optimal-policy-appendix}
  Suppose $\E[\ATE_3^2]<\infty$, $\var(\ATE_3)>0$, and the uniform boundary anti-concentration and ordinary sign-separation conditions \Cref{eq:uniform-boundary-anti-concentration-large-budget-appendix,eq:ordinary-sign-separation-large-budget-appendix} hold.
  Then:
  \begin{enumerate}
    \item[\rm (i)] $V^\star(b)\le V_{\oracle}$ for all $b>0$.
    \item[\rm (ii)] $V^\star(b)\to V_{\oracle}$ as $b\to\infty$.
    \item[\rm (iii)] The optimal investment diverges,
      \[
        \inf_{\bm s_1^\star\in\mathcal S_1^\star(b)} \bm c^\top\bm s_1^\star
        \to\infty
        \qquad\text{as }b\to\infty.
      \]
    \item[\rm (iv)] The set of unit-investment normalized optimal designs converges to the unit-investment representative design,
      \[
        \sup_{\bm s_1^\star\in\mathcal S_1^\star(b)}
        \left\|
        \frac{\bm s_1^\star}{\bm c^\top\bm s_1^\star}
        -
        \overline{\bm s}_1^{\,\infty}
        \right\|
        \to0
        \qquad\text{as }b\to\infty.
      \]
  \end{enumerate}
\end{linkedresult}

\begin{proof}
  Apply \Cref{thm:large-budget-weighted-target-impact-appendix} with $Z=\Test_2$.
  By \Cref{lem:test-two-satisfies-weighted-conditions-appendix}, $Z$ is non-adaptive, $0\le Z\le1$, and $m_Z(\bm\tau)=p_2(\ATE_2)>0$ almost surely.
  Since $\var(\ATE_3)>0$ and the ordinary sign-separation condition holds, \Cref{prop:positive-conditional-weight-appendix} verifies the weighted non-degeneracy and weighted sign-separation conditions.
  For $Z=\Test_2$, the generic value, oracle, noise-free value, and optimal solution set are exactly $V$, $V_{\oracle}$, $V_{\noisefree}$, and $\mathcal S_1^\star(b)$.
\end{proof}

\begin{proposition}[Density implies large-budget regularity]
  \label{prop:density-large-budget-regularity-appendix}
  Suppose the prior distribution of $\bm\tau$ is absolutely continuous with respect to Lebesgue measure on $\R^D$, with density $f$.
  Assume there exists $r>0$ such that
  \[
    f(x)>0
    \qquad
    \text{for Lebesgue-a.e. }x\in B_r\defeq\{x\in\R^D:\|x\|_2<r\}.
  \]
  Then $\ATE_3$ is non-atomic.
  In particular, if $\E[\ATE_3^2]<\infty$, then $\var(\ATE_3)>0$.
  Moreover, the uniform boundary anti-concentration condition \Cref{eq:uniform-boundary-anti-concentration-large-budget-appendix} and the ordinary sign-separation condition \Cref{eq:ordinary-sign-separation-large-budget-appendix} hold.
  Consequently, under the finite-second-moment assumption in the body theorem, the hypotheses of \linkedCref{thm:large-budget-limiting-optimal-policy-appendix} are satisfied.
\end{proposition}

\begin{proof}
  \emph{Step 1: Non-atomicity and positive variance of $\ATE_3$.}
  The target sample is nonzero, i.e. $\bm s_3\neq\bm0$, and for any $a\in\R$,
  \[
    \left\{\bm{\tau}\in\R^D:\ATE_3=a\right\}
    =
    \left\{\bm{\tau}\in\R^D:\bm s_3^\top \bm{\tau}=a\,\bm1^\top\bm s_3\right\}.
  \]
  is a hyperplane and therefore has Lebesgue measure zero.
  Absolute continuity of the law of $\bm\tau$ implies $\P(\ATE_3=a)=0$ for every $a\in\R$.
  Thus $\ATE_3$ has no atoms.
  If $\E[\ATE_3^2]<\infty$ and $\var(\ATE_3)=0$, then $\ATE_3$ would be almost surely equal to the constant $\E[\ATE_3]$, contradicting non-atomicity.
  Hence $\var(\ATE_3)>0$.

  \emph{Step 2: Uniform boundary anti-concentration condition.}
  For $\overline{\bm s}_1\in\overline{\mathcal{S}}$, define
  \[
    \bm u(\overline{\bm s}_1)
    \defeq
    \frac{\overline{\bm s}_1}{\sqrt{\bm 1^\top \overline{\bm s}_1}},
    \qquad
    \xi(\overline{\bm s}_1,\bm\tau)
    =
    \bm u(\overline{\bm s}_1)^\top \bm\tau.
  \]
  Because $\overline{\mathcal{S}}$ is compact and
  \[
    \bm 1^\top \overline{\bm s}_1
    \ge
    \frac{1}{\|\bm c\|_\infty}
    \qquad
    \forall \overline{\bm s}_1\in\overline{\mathcal{S}},
  \]
  the image
  \[
    K
    \defeq
    \bigl\{\bm u(\overline{\bm s}_1):\overline{\bm s}_1\in\overline{\mathcal{S}}\bigr\}
    \subset \R^D\setminus\{\bm 0\}
  \]
  is compact.

  We claim that
  \[
    \sup_{\overline{\bm s}_1\in\overline{\mathcal{S}}}
    \P\!\bigl(|\xi(\overline{\bm s}_1,\bm\tau)|\le \delta\bigr)
    =
    \sup_{\bm u\in K}\P\!\bigl(|\bm u^\top\bm\tau|\le \delta\bigr)
    \longrightarrow 0
    \qquad\text{as }\delta\downarrow 0.
  \]
  Suppose not.
  Then there exist $\epsilon_0>0$, a sequence $\delta_n\downarrow 0$, and $\bm u_n\in K$ such that
  \[
    \P\!\bigl(|\bm u_n^\top\bm\tau|\le \delta_n\bigr)\ge \epsilon_0
    \qquad
    \forall n.
  \]
  By compactness of $K$, after passing to a subsequence we may assume $\bm u_n\to \bm u\in K$.

  Fix $x\in\R^D$ such that $\bm u^\top x\neq 0$.
  Since $\bm u_n^\top x\to \bm u^\top x$ and $\delta_n\downarrow 0$, we have
  \[
    \indic{|\bm u_n^\top x|\le \delta_n}\longrightarrow 0.
  \]
  The exceptional set $\{x:\bm u^\top x=0\}$ is a hyperplane and therefore has Lebesgue measure zero.
  Since the law of $\bm\tau$ is absolutely continuous, it follows that
  \[
    \indic{|\bm u_n^\top\bm\tau|\le \delta_n}\longrightarrow 0
    \qquad\text{a.s.}
  \]
  Moreover, these indicators are bounded by $1$, so dominated convergence gives
  \[
    \P\!\bigl(|\bm u_n^\top\bm\tau|\le \delta_n\bigr)
    =
    \E\!\left[\indic{|\bm u_n^\top\bm\tau|\le \delta_n}\right]
    \longrightarrow 0,
  \]
  contradicting the lower bound by $\epsilon_0$.
  This proves the uniform boundary anti-concentration condition.

  \emph{Step 3: Sign separation.}
  Fix $\overline{\bm s}_1\in\overline{\mathcal{S}}$ not proportional to $\bm s_3$.
  Let $\bm u\defeq \overline{\bm s}_1$.
  Since $\bm u$ is not proportional to $\bm s_3$, there exists $x_0\in\R^D$ such that
  \[
    \bm u^\top x_0 = 0,
    \qquad
    \bm s_3^\top x_0\neq 0.
  \]
  Replacing $x_0$ by $-x_0$ if necessary, assume $\bm s_3^\top x_0<0$.
  Then for any $\eta>0$,
  \[
    \bm u^\top(x_0+\eta \bm u)=\eta\|\bm u\|_2^2>0,
  \]
  and, for all sufficiently small $\eta>0$,
  \[
    \bm s_3^\top(x_0+\eta \bm u)
    =
    \bm s_3^\top x_0+\eta\,\bm s_3^\top \bm u
    <0.
  \]
  Hence the set
  \[
    A
    \defeq
    \{x\in\R^D:(\overline{\bm s}_1^\top x)(\bm s_3^\top x)<0\}
  \]
  is nonempty.
  It is also open, and it is a cone: if $x\in A$ and $t>0$, then $tx\in A$.
  Therefore $A\cap B_r$ is a nonempty open subset of $B_r$.

  Since $f>0$ for Lebesgue-a.e. $x\in B_r$, we have $f>0$ for Lebesgue-a.e. $x\in A\cap B_r$.
  Because $A\cap B_r$ is nonempty and open, it has strictly positive Lebesgue measure, and thus
  \[
    \P\!\bigl((\overline{\bm s}_1^\top\bm\tau)(\bm s_3^\top\bm\tau)<0\bigr)
    \ge
    \P(\bm\tau\in A\cap B_r)
    =
    \int_{A\cap B_r} f(x)\,dx
    >
    0.
  \]
  This proves sign separation.
\end{proof}
